\documentclass[11pt]{article}
\usepackage[T1]{fontenc}
\usepackage{textcomp}
\usepackage[margin=0.75in]{geometry}
\usepackage{amsmath,amssymb,amsthm}
\usepackage{array,booktabs}
\usepackage{hyperref}
\setlength{\parindent}{0pt}
\newtheorem{theorem}{Theorem}
\theoremstyle{definition}
\newtheorem{definition}{Definition}
\title{Flow Proof Artifact: Group-inv-inv}
\author{Generated by \texttt{flow doc proof}}
\date{Flow Proof Book}
\begin{document}
\maketitle
The inverse of an inverse returns the original element.\\[0.5em]
\textbf{Source.} dummit-foote — *Abstract Algebra*, §1.1\\[0.5em]
\bigskip
\noindent{\large \textbf{Derived fact 1.} \textit{inv(inv(g)) = g} \hfill Group \cdot inverse \cdot double inverse returns the element}\par
\smallskip\noindent\textit{Coordinate: Group \cdot inverse \cdot double inverse returns the element}\par
\smallskip\noindent\textit{Source: dummit-foote}\par
\smallskip\noindent\textit{Built on: left inverse recovers the identity, for inverse on Group, inverses are unique, for inverse on Group}\par
\medskip
\noindent\textbf{Goal.} inv(inv(g)) = g\par
\noindent$\displaystyle \forall g \in Group\quad inv(inv(g)) = g$\par
\medskip
\noindent
\begin{tabular}{@{} >{\raggedright\arraybackslash}p{0.06\textwidth} >{\raggedright\arraybackslash}p{0.47\textwidth} >{\raggedleft\arraybackslash}p{0.41\textwidth} @{}}
\textbf{\#} & \textbf{Proof} & \textbf{Mathematics} \\
\midrule
\textbf{1.} & We prove that double inverse returns the element for inverse on Group. & \\[0.45em]
\textbf{2.} & We invoke the definitional clause governing inverse on Group: left inverse recovers the identity, for inverse on Group (instantiated for inv(g)). & \\[0.45em]
\textbf{3.} & We invoke the derived fact governing inverse on Group: inverses are unique, for inverse on Group (instantiated for g, inv(inv(g)), g). & \\[0.45em]
\textbf{4.} & From step 2 and step 3, this implies inv(inv(g)) equals g. Hence proven. & $\displaystyle inv(inv(g)) = g$ \\[0.45em]
\bottomrule
\end{tabular}
\label{thm:Group-inverse-double_inverse_returns_the_element}
\par\medskip\hrule
\end{document}
